\documentclass{cv}

\usepackage{fontawesome5}
\usepackage{multicol}

\cvname{Dylan Green}
\cvemail{dylan.green@eng.ox.ac.uk}
\cvlinkedin{d-green}
\cvwebppage{https://d-green1958.github.io/}
\cvgithub{d-green1958}
\cvorcid{0009-0008-2431-0587}
\cvgooglescholar{https://scholar.google.com/citations?user=v8EhQi4AAAAJ}
\cvphoneno{+44 7776 796 498}
\cvtitle{PhD Researcher | Department of Engineering Science}
\cvaddress{Flat 18 Osney Mews OX2 0PF}

\cvbio{PhD student at the University of Oxford researching floating wind turbine wake dynamics using high-performance computing. Interests include computational fluid dynamics and large-scale simulation, with a focus on applications to renewable energy engineering.}

\begin{document}
\makecvheader

\cvsection{Education}
\cvitem{PhD (DPhil) Wind \& Marine Energy Systems \& Structures}{University of Oxford}{Oxford, Uk}{2023--Present}
\begin{itemize}
	\item Investigated the impact of floating platform motions on the floating wind turbine wakes using the Actuator Line Model (ALM).
	\item Developed and extended the in‑house OpenFOAM‑based ALM code.
	\item Performed high fidelity large eddy simulations using OpenFOAM on high‑performance computing (HPC) clusters to resolve transient wake behaviour.
	\item Validated numerical model against open‑source experimental results.
	\item Extensive use of ARCHER2 and ARC HPCs.
	\item Used Python and Paraview for data visualisation.
\end{itemize}
\cvitemskip

\cvitem{EPSRC Centre for Doctoral Training}{University of Strathclyde}{Glasgow, Uk}{2023--2024}
\begin{itemize}
	\item Short intensive courses providing a holistic understanding of offshore renewable energy engineering.
\end{itemize}
\cvitemskip

\cvitem{MMathPhys Master of Mathematics \& Physics with Honours}{University of Manchester}{Manchester, Uk}{2019--2023}
\begin{itemize}
	\item[] \textbf{Grade: First‑class honours (82\%)}
	\item Took a mix of applied and therotical courses focusing on continuum mechanics, numerical methods and scientific computing.
	\item MMath Project: \textit{Finite Element Solutions to the Helmholtz Equation}. Developed a C++ finite element solver for the Helmholtz equation. Assessed boundary conditions for wave scattering problems, including perfectly matched layers, Dirichlet‑to‑Neumann maps, and approximate radiation conditions
	\item MPhys Project: \textit{Quantifying Correlation Between Geometric Uncertainties in Rectum Contouring and The Resulting Patient Toxicity}. Applied a deep learning algorithm to standardize rectal contouring. Performed statistical analysis to determine the correlation between patient outcomes and contour deviations. Resulted in publication.
\end{itemize}


\cvsection{Experience}
\cvitem{Conference Organiser}{Wind \& Marine Systems \& Structures Centre for Doctoral Training}{Oxford, Uk}{2026--Present}
\begin{itemize}
	\item Organiser of student-led Future Wind \& Marine Conference hosted in Oxford in April 2027.
	\item Responsible for fundraising and bookkeeping of event.
\end{itemize}
\cvitemskip


\cvitem{Lab Demonstrator}{Department of Engineering Science, University of Oxford}{Oxford, Uk}{2025--Present}
\begin{itemize}
	\item Supervised \& supported several groups of 24 students in the material testing and computational labs.
	\item Assisted in experimental design, data interpretation and provided one‑on‑one support for algorithm design and debugging.
	\item Responsible for marking through interview style assessment
\end{itemize}
\cvitemskip

\cvitem{Research Intern}{Manchester Centre for Nonlinear Dynamics, University of Manchester}{Manchester, Uk}{2022}
\begin{itemize}
	\item Project: \textit{Biometric Modelling of Retinal Injections} - Experimental work aimed at identifying nonlinear dynamics in retinal injections focusing on viscous fingering instabilities.
	\item Designed and executed fluid injection experiments, analysed and presented results to multi-disciplinary research team.
\end{itemize}
\cvitemskip

\cvitem{Work Experience: Low Frequency Demand Disconnection Team}{National Grid}{Warwick, Uk}{2019}
\begin{itemize}
	\item Worked as part of Low Frequency Demand Disconnection team wherein I utilised low‑order models to optimise demand disconnection procedure during the event of generation failure. Produced report summarising effects of sudden generation failure on system stability.
\end{itemize}


\cvsection{Skills}
{
%\vspace{-1em}
%\begin{multicols}{2}
\begin{itemize}[topsep=0pt, itemsep=1pt]
	\item[\cvbodysz\faCode] \cvbodysz \textbf{Programming} Python | \texttt{C++} | \texttt{C} | MatLab | Linux
	\item[\cvbodysz\faHammer] \cvbodysz \textbf{Build \& Job Automation} {bash} | {SLURM} | {Make} | {CMake}
	\item[\cvbodysz\faSlidersH] \cvbodysz \textbf{Data Analysis} Numpy | Paraview | \texttt{VTK}
	\item[\cvbodysz\faBolt] \cvbodysz \textbf{High-Performance Computing} {CUDA} | {HIP} | {MPI} | {MP}
	\item[\cvbodysz\faTextHeight] \cvbodysz \textbf{Type Setting} \LaTeX | Beamer | Word
	\item[\cvbodysz\faGit*] \cvbodysz \textbf{Version Control} {Git} | GitHub
	\item[\cvbodysz\faTint] \cvbodysz \textbf{CFD / Numerical Methods} OpenFoam | FEA
\end{itemize}
%\end{multicols}
}

\newpage
\cvsection{Courses \& Certificates}
\begin{itemize}
	\item \textbf{NVidia} - Fundamentals of Accelerated Copmuting With Modern CUDA C++
	\item \textbf{Zero Institute} - Climate Tech Hackathon First Place
	\item \textbf{University of Oxford} - Course on CUDA Programming on NVIDIA GPUs
	\item \textbf{ARCHER2 National Super Computer} - GPU programming with kernels
\end{itemize}

\cvsection{Publications \& Presentations}
\begin{itemize}
	\item \textit{Support Structure Impacts on Floating Offshore Wind Turbine Wake Aerodynamics} - Paper \& Poster, The Science of Making Torque from Wind, Bruges 2026.
	\item \textit{Energy \& Momentum Transport in The Wakes of Floating Turbines} - Presentation, Recent Advances in Turbulent Wind Farm Dynamics Colloquium, London 2026. 
	\item \textit{Sensitivity of the Floating Actuator Line Method to Prescribed Motions} - Presentation, 21st European Academy of Wind Energy PhD Seminar PhD Seminar, Athens 2025.
	\item \textit{Development and Validation of an Actuator Line Method for Floating Turbines} - Paper \& Presentation, XI Conference on Computational Methods in Marine Engineering, Edinburgh 2025.
	\item \textit{New rectum dose surface mapping methodology to identify rectal subregions associated with toxicities following prostate cancer radiotherapy} - paper, Physics and Imaging in Radiation Oncology 2025.
\end{itemize}


\end{document}
